% TOP_PROBLEM: {"id": 30007067, "title": "Two-step linear images of special quasi-projective groups", "category_id": 5, "category": {"id": 5, "name": "algebraic_geometry", "display_name": "Algebraic Geometry", "description": "Geometric objects defined by polynomial equations.", "slug": "algebraic-geometry", "order_index": 5, "created_at": "2026-07-31T15:26:25.671Z"}, "status": "open"}
% FIRST_POSTED: 2026-10-06
% ENTRY_KIND: statement_only
% !TeX program = lualatex
\documentclass[11pt]{article}
\usepackage[margin=25mm]{geometry}
\usepackage{fontspec}
\setmainfont{Latin Modern Roman}
\usepackage{amsmath,amssymb,mathtools}
\usepackage{unicode-math}
\setmathfont{Latin Modern Math}
\usepackage{xurl,hyperref}
\hypersetup{colorlinks=true,urlcolor=blue}
\setcounter{secnumdepth}{0}
\setlength{\parindent}{0pt}
\setlength{\parskip}{5pt}
\setlength{\emergencystretch}{3em}
\newcommand{\R}{\mathbb R}
\newcommand{\N}{\mathbb N}
\newcommand{\Z}{\mathbb Z}
\newcommand{\Q}{\mathbb Q}
\newcommand{\C}{\mathbb C}
\newcommand{\D}{\mathbb D}
\newcommand{\F}{\mathbb F}
\newcommand{\sref}[2]{\hyperlink{source:#1:#2}{[S#2]}}
\newcommand{\cataloguescope}[1]{#1}
\begin{document}
\section*{Two-step linear images of special quasi-projective groups}\label{30007067}
\subsection{Definitions and mathematical statement}
\paragraph{Statement recorded in the supplied source.}
For every connected smooth special complex quasi-projective variety V, every $d\geq 1$ and every representation $\rho :\pi _{1}(V)\to GL_d(C)$, must its image have a finite-index subgroup nilpotent of class at most two? The assertion concerns representation images, not virtual nilpotence of $\pi _{1}(V)$ itself.
\subsection{Short English statement}
For every connected smooth special complex quasi-projective variety V, every $d\geq 1$ and every representation $\rho :\pi _{1}(V)\to GL_d(C)$, must its image have a finite-index subgroup nilpotent of class at most two? The assertion concerns representation images, not virtual nilpotence of $\pi _{1}(V)$ itself.
\subsection{Sources}
\begin{itemize}
\item[S1] Original question source or mathematical context.
\url{http://aimpl.org/amenablediscrete/5/}.
\item[S2] Original question source or mathematical context.
\url{https://doi.org/10.4171/owr/2014/09}.
\item[S3] Original question source or mathematical context.
\url{https://link.springer.com/article/10.1007/s10231-024-01501-8}.
\item[S4] OpenAI, manuscript citation page. Source of the formulation motivating this local problem.
\url{https://github.com/openai/math/blob/adc7f1241b42e322a6451854ab7e4b4c146bf78a/preprints/Two-step-monodromy-of-special-quasi-projective-varieties-September-24-2026/README.md}.
\end{itemize}
\end{document}
